%% Chapter 1 — Introduction
\section{Introduction}
\label{sec:introduction}

Quantum computing holds the promise of solving classically intractable problems
in optimisation, simulation, and cryptography.  Among the leading hardware
platforms, superconducting qubits stand out for their relatively fast gate times,
high coherence, and compatibility with existing microwave engineering
infrastructure~\cite{Krantz2019}.  A central challenge that must be overcome on
the road to fault-tolerant operation is \emph{qubit state measurement}: reading
out the quantum information stored in a qubit quickly, accurately, and with
minimal back-action on unmeasured qubits.

\subsection{Motivation}

Classical error correction requires measurements that are not only high-fidelity
but also \emph{fast} relative to qubit coherence times, and — crucially for
quantum error correction — it requires \emph{mid-circuit measurements}: the ability to measure ancilla qubits repeatedly during a quantum algorithm, use the results to extract error syndromes in real time, and continue the computation on the data qubits.  This capability demands a control electronics stack capable of making a state discrimination decision on-chip in microseconds — the timescale of the readout integration window itself — rather than the milliseconds typical of a software round-trip to a host CPU.

Field-Programmable Gate Arrays (FPGAs) offer a compelling solution.  Their
deterministic, sub-microsecond timing enables real-time signal processing close
to the qubit hardware.  Commercial control platforms such as Qblox, Zurich
Instruments, and Quantum Machines already exploit this, providing on-FPGA state
discrimination within closed, proprietary firmware stacks.  However, closed
firmware precludes custom signal processing pipelines: researchers cannot inject
bespoke algorithms — including machine-learning classifiers — into the real-time
decision path.  The Quantum Instrumentation Control Kit
(\QICK)~\cite{Stefanazzi2022} addresses this gap as an open-source FPGA firmware
and software framework that turns Xilinx Radio-Frequency System-on-Chip (RFSoC)
boards into complete qubit controllers, combining fast digital-to-analogue
converters (DACs), analogue-to-digital converters (ADCs), and a tightly
integrated soft-processor into a single compact platform.  Full access to the
programmable logic fabric makes it possible to define the entire readout and
discrimination pipeline, including the integration of custom ML models, in a way
that closed commercial firmware does not permit.

\subsection{Goals and Contributions}

Despite the availability of \QICK{}, assembling a \emph{general-purpose} readout
pipeline from an uncharacterised board to a working single-shot state
discriminator requires integrating many layers: physical wiring, FPGA
configuration, phase calibration, signal optimisation, and machine-learning
post-processing.  The lack of a unified reference implementation makes this
process time-consuming and error-prone, especially for multi-qubit scenarios.

This thesis makes the following contributions:

\begin{enumerate}
      \item \textbf{Full-stack integration} of the \QICK{} framework on a Xilinx
            ZCU216 RFSoC board with an XM655 analogue front-end module, including
            remote operation via a persistent Pyro4 server.

      \item \textbf{Phase-coherent readout calibration}: a systematic procedure for
            measuring and compensating the random DAC–ADC phase offset that arises at
            every board power-on, enabling reproducible IQ measurements across
            sessions.

      \item \textbf{Readout data-path design}: a characterisation of the trade-offs
            between decimated and accumulated readout modes and a complete resonator
            spectroscopy workflow suitable for both single-qubit and multiplexed
            measurements.

      \item \textbf{Real-time feedback primitives}: demonstration of the
            \texttt{condj} conditional-branch primitive on the tProcessor in loopback,
            and design of an active qubit reset sequence combining dispersive readout,
            thresholding, and a calibrated $\pi$ pulse — establishing the
            infrastructure needed for mid-circuit measurement once the board is
            connected to a qubit chip.

      \item \textbf{ML-assisted state discrimination}: offline design and evaluation
            of a lightweight ML discrimination pipeline on existing experimental SSRO
            data, establishing that LDA is near-optimal for integrated IQ inputs, and
            demonstrating via synthetic raw traces that a 1D CNN recovers
            $T_1$-relaxation shots irrecoverable by integration-based methods.  A
            concrete path toward FPGA deployment is outlined.
\end{enumerate}

\subsection{Thesis Outline}

Section~\ref{sec:background} reviews the relevant theory of superconducting
qubits, circuit quantum electrodynamics (cQED), and dispersive readout.
Section~\ref{sec:hardware} describes the hardware platform.
Section~\ref{sec:integration} covers system integration and calibration.
Section~\ref{sec:readout} presents the readout architecture and measurement
workflows.  Section~\ref{sec:feedback} describes real-time feedback and
mid-circuit measurement.  Section~\ref{sec:ml} evaluates machine-learning-based
state discrimination.  Section~\ref{sec:conclusions} summarises the results and
outlines future directions.
